\begin{table}[htbp]
\centering
\footnotesize
\caption{Paired comparison of simulated control outcomes across all 30 fitted participants.}
\label{tab:control-cohort}
\begin{tabular}{p{.22\linewidth}p{.14\linewidth}p{.14\linewidth}p{.24\linewidth}p{.13\linewidth}}
\hline
Outcome & No control & Optimized & Change [95\% CI] & Better/same/worse \\
\hline
Time in 70-110 mg/dL (\%) & 66.75 (38.02) & 94.13 (12.48) & +27.37 [+16.56, +39.45] & 18/11/1 \\
Time above 110 mg/dL (\%) & 32.83 (38.33) & 4.79 (12.51) & -28.04 [-39.87, -17.28] & 19/11/0 \\
Time below 70 mg/dL (\%) & 0.42 (2.30) & 1.08 (3.16) & +0.66 [+0.00, +1.59] & 0/27/3 \\
Peak glucose (mg/dL) & 123.70 (33.74) & 102.60 (17.56) & -21.09 [-27.93, -14.77] & 26/4/0 \\
Paper objective J (model objective units) & 902.88 (2909.14) & -322.50 (451.05) & -1225.39 [-2216.09, -465.13] & 26/4/0 \\
Integrated control intensity (model intensity x h) & 0.00 (0.00) & 4.30 (4.83) & +4.30 [+2.71, +6.09] & -- \\
\hline
\end{tabular}
\par\smallskip
\begin{minipage}{\textwidth}
\scriptsize
Values are mean (SD). Change = optimized minus no control; time changes are percentage points. Paired percentile bootstrap CIs use 10,000 resamples. One median-meal simulation per fitted participant; conditional on fitted model parameters. Control is uncalibrated model intensity. Full definitions and scope are in the analysis report.
\end{minipage}
\end{table}
